% !TEX root = 08.tex
\documentclass[a4paper,10pt]{article}

\usepackage{pdfsync}
\usepackage[utf8]{inputenc}
\usepackage{microtype}
\usepackage{mathtools}
\usepackage{amsthm,amsfonts,amsmath,amssymb}
\usepackage{mathabx}
% \usepackage{MnSymbol}
\usepackage{hyperref}
\usepackage{bbold}
\usepackage[textwidth=2cm,textsize=small]{todonotes}
\usepackage{subcaption}
\usepackage{enumerate}
\usepackage{xspace}
\usepackage{stmaryrd}
\usepackage{verbatim}
\usepackage{cleveref}

\hypersetup{
    colorlinks,
    linkcolor={red!50!black},
    citecolor={blue!50!black},
    urlcolor={blue!80!black}
}

\author{}

\newcommand{\ignore}[1]{}

\newcommand{\st}{s.t.}
\newcommand{\ie}{i.e.}
\newcommand{\eg}{e.g.}

\newcommand{\sep}{\;|\;}
\renewcommand{\rule}[2]{\dfrac{#1}{#2}}
\newcommand{\set}[1]{\left\{#1\right\}}
\newcommand{\setof}[2]{\set{#1 \;\middle|\; #2}}
\newcommand{\tuple}[1]{\left( #1 \right)}
\newcommand{\pair}[2]{\tuple {#1, #2}}
\newcommand{\closure}[3]{\tuple {#1, #2, #3}}
\newcommand{\sem}[1]{\llbracket{#1}\rrbracket}
\newcommand{\powerset}[1]{\mathcal P(#1)}
\newcommand{\map}[2]{\mathsf{map}\; #1\; #2}
\newcommand{\filter}[2]{\mathsf{filter}\; #1\; #2}
\newcommand{\floor}[1]{\left\lfloor #1 \right\rfloor}
\newcommand{\PreRel}[2]{\mathsf{Pre}_{#1}(#2)}
\newcommand{\Pre}[1]{\PreRel {} {#1}}
\newcommand{\PreStar}[1]{\mathsf{Pre}^*(#1)}
\newcommand{\Post}[1]{\mathsf{Post}(#1)}
\newcommand{\PostStar}[1]{\mathsf{Post^*}(#1)}
\newcommand{\Conf}{\mathsf{Conf}}
\newcommand{\upwardclosure}[1]{#1\uparrow}
\newcommand{\card}[1]{\left| #1 \right|}
\newcommand{\lang}[1]{L(#1)}
\newcommand{\goesto}[1]{\xrightarrow{#1}}
\newcommand{\e}{\varepsilon}

\newcommand{\N}{\mathbb N}
\newcommand{\Z}{\mathbb Z}
\newcommand{\Q}{\mathbb Q}

% LTL
\newcommand{\true}{\mathbf{true}}
\newcommand{\false}{\mathbf{false}}
\newcommand{\X}{\mathbf X}
\newcommand{\U}{\mathbf U}
\newcommand{\Release}{\mathbf R}
\newcommand{\F}{\mathbf F}
\newcommand{\G}{\mathbf G}

% CTL
\newcommand{\EX}{\mathbf{EX}}
\newcommand{\AX}{\mathbf{AX}}
\newcommand{\EF}{\mathbf{EF}}
\newcommand{\AF}{\mathbf{AF}}
\newcommand{\EG}{\mathbf{EG}}
\newcommand{\AG}{\mathbf{AG}}
\newcommand{\EU}[2]{\mathbf E(#1 \U #2)}
\newcommand{\AU}[2]{\mathbf A(#1 \U #2)}
\newcommand{\ER}[2]{\mathbf E(#1 \Release #2)}
\newcommand{\AR}[2]{\mathbf A(#1 \Release #2)}

\newcommand{\NL}{\textsf{NL}}
\newcommand{\PTIME}{\textsf{PTIME}}
\newcommand{\PSPACE}{\textsf{PSPACE}}

\makeatletter

\def\dasharrowfill@#1#2#3#4{%
        $\m@th
        \thickmuskip0mu
        \medmuskip\thickmuskip
        \thinmuskip\thickmuskip
        \relax
        #4#1\mkern2mu
        \xleaders\hbox{$#4\mkern2mu#2\mkern2mu$}\hfill
        \mkern2mu
        #3$%
}

\def\dashleftarrowfill@{\dasharrowfill@\leftarrow\relbar\relbar}
\def\dashrightarrowfill@{\dasharrowfill@\relbar\relbar\rightarrow}
\def\dashleftrightarrowfill@{\dasharrowfill@\leftarrow\relbar\rightarrow}
\def\dashLeftarrowfill@{\dasharrowfill@\Leftarrow\Relbar\Relbar}
\def\dashRightarrowfill@{\dasharrowfill@\Relbar\Relbar\Rightarrow}
\def\dashLeftrightarrowfill@{\dasharrowfill@\Leftarrow\Relbar\Rightarrow}

\providecommand*\xdashleftarrow[2][]{%
  \ext@arrow 0055{\dashleftarrowfill@}{#1}{#2}}
\providecommand*\xdashrightarrow[2][]{%
  \ext@arrow 0055{\dashrightarrowfill@}{#1}{#2}}
\providecommand*\xdashleftrightarrow[2][]{%
  \ext@arrow 0055{\dashleftrightarrowfill@}{#1}{#2}}
\providecommand*\xdashLeftarrow[2][]{%
  \ext@arrow 0055{\dashLeftarrowfill@}{#1}{#2}}
\providecommand*\xdashRightarrow[2][]{%
  \ext@arrow 0055{\dashRightarrowfill@}{#1}{#2}}
\providecommand*\xdashLeftrightarrow[2][]{%
  \ext@arrow 0055{\dashLeftrightarrowfill@}{#1}{#2}}

\def\slashedarrowfill@#1#2#3#4#5{%
$\m@th\thickmuskip0mu\medmuskip\thickmuskip\thinmuskip\thickmuskip
  \relax#5#1\mkern-7mu%
  \cleaders\hbox{$#5\mkern-2mu#2\mkern-2mu$}\hfill
  \mathclap{#3}\mathclap{#2}%
  \cleaders\hbox{$#5\mkern-2mu#2\mkern-2mu$}\hfill
  \mkern-7mu#4$%
}

\def\rightslashedarrowfill@{%
  \slashedarrowfill@\relbar\relbar{\raisebox{1.2pt}{$\scriptscriptstyle\diagup$}}\rightarrow}
\newcommand\xslashedrightarrow[2][]{%
  \ext@arrow 0055{\rightslashedarrowfill@}{#1}{#2}}
\makeatother

\newtheorem{exercise}{Exercise}

\newenvironment{solution}{\begin{proof}[\protect{\textnormal{\emph{Solution.}}}]}{\end{proof}}

\let\oldqedhere\qedhere

\newif\ifstartedinmathmode
\renewcommand*{\qedhere}{%
  \relax\ifmmode\startedinmathmodetrue\else\startedinmathmodefalse\fi%
  \ifstartedinmathmode\tag*{\protect\oldqedhere}\else\ifinlist{\hfill\oldqedhere}{\oldqedhere}\fi%
}

\crefname{equation}{}{}
\crefname{exercise}{Exercise}{Exercises}


\title{Formal verification - Tutorial 08}
\date{11/05/2026}

\begin{document}

\maketitle

\section*{Büchi's Büchi automata complementation}

Let $D = \set{\bot, 0, 1}$ be a set with three elements.
Intuitively, $\bot$ means no run, $0$ means a run,
and $1$ means a run visiting an accepting transition.
%
In the following, fix a set of states $Q$.
%
A \emph{graph} is a function $g \in D^{Q \times Q}$.
%
A Büchi automaton is a tuple $A = \tuple{\Sigma, Q, I, \delta}$ where
$\Sigma$ is the input alphabet,
$Q$ is the set of states, $I \subseteq Q$ is the set of initial states,
and the transition relation is represented by a family of graphs
$\delta : \Sigma \to D^{Q \times Q}$,
one graph $\delta_a$ for each input symbol $a \in \Sigma$.
%
The language $\lang A$ recognised by $A$ is defined as expected.

We study how to build the complement of $A$,
following Büchi's original Ramsey-based construction.
%
For two states $p, q \in Q$ and an input word $w \in \Sigma^*$, write
%
\begin{itemize}
    \item $p \xrightarrow w q$ if there is a run from $p$ to $q$ on $w$,
    \item $p \xdashrightarrow w q$ if there is a run from $p$ to $q$ on $w$ visiting an accepting transition, and
    \item $p \xslashedrightarrow w q$ if there is no run from $p$ to $q$ on $w$.
\end{itemize}
%
An \emph{edge} is a tuple $e = \tuple {p, x, q}$ where $p, q \in Q$ and $x \in D$.
The \emph{language} (or \emph{semantics}) of an edge is defined as follows.
%
\begin{align}
    \lang {p, 1, q}
        &= \setof{w \in \Sigma^*}{p \xdashrightarrow w q}, \\
    \lang {p, 0, q}
        &= \setof{w \in \Sigma^*}{p \xrightarrow w q} \setminus \lang {p, 1, q}, \\
    \lang {p, \bot, q}
        &= \setof{w \in \Sigma^*}{p \xslashedrightarrow w q}
        = \Sigma^* \setminus (\lang {p, 0, q} \cup \lang {p, 1, q}).
\end{align}
%
The language of a graph $g$ is the intersection of the languages of its edges,
%
\begin{align}
    \lang g = \bigcap_{e \in g} \lang e = \bigcap_{p, q \in Q} \lang {p, g(p, q), q}.
\end{align}
%
Finally, the language of two graphs $g, h$ is
%
\begin{align}
    \lang {g, h} = \lang g \cdot \lang h^\omega.
\end{align}
%
We use languages $\lang {g, h}$ to construct the complement of $\lang A$.

\begin{exercise}
    Show that for every finite word $w \in \Sigma^*$
    there exists a graph $g_w$, the \emph{characteristic graph} associated to $w$, \st~$w \in \lang {g_w}$.
\end{exercise}
\begin{solution}
    Let $g_w$ be the graph defined as follows.
    For every pair of states $p, q \in Q$,
    we have $g_w(p, q) = 1$ if $p \xdashrightarrow w q$,
    $g_w(p, q) = 0$ if $p \xrightarrow w q$ but $p \xslashedrightarrow w q$, and
    $g_w(p, q) = \bot$ if $p \xslashedrightarrow w q$.
    %
    By definition, $w$ is in the language of every edge of $g_w$,
    and thus $w \in \lang {g_w}$, as required.
\end{solution}

\begin{exercise}
    Show that $\Sigma^\omega = \bigcup_{g, h} \lang {g, h}$,
    where the union is taken over all pairs of graphs $g, h$.
\end{exercise}
\begin{solution}
    We apply the infinite Ramsey's theorem.
    Consider an arbitrary infinite word $w = a_0 a_1 \cdots \in \Sigma^\omega$.
    %
    Consider the infinite graph with nodes $\N$
    and edges $\tuple {i, j}$ for $i < j$ coloured by $g_{a_i a_{i+1} \cdots a_{j-1}}$.
    %
    Since there are only finitely many colours,
    by the infinite Ramsey's theorem the graph as an infinite monochromatic subgraph.
    %
    In other words, there exist an infinite set $I = \set{i_0 < i_1 < \cdots} \subseteq \N$ and a graph $h$
    \st~for every $i_j$ from $I$ we have $g_{a_{i_j} a_{i_j+1} \cdots a_{i_{j+1}-1}} = h$.
    %
    We can accordingly partition the word $w$ into segments
    %
    \begin{align*}
        w = w_0 w_1 w_2 \cdots
    \end{align*}
    %
    \st~$g := g_{w_0}$ and $h = g_{w_1} = g_{w_2} = \cdots$.
    %
    % (In particular, $h$ is idempotent $h = h; h$.)
    %
    It follows that $w \in \lang {g, h}$, as required.
\end{solution}

\begin{exercise}
    For every pair of graphs $g, h$,
    show that either $\lang {g, h} \subseteq \lang A$ or $\lang {g, h} \cap \lang A = \emptyset$.
\end{exercise}
%
\begin{solution}
    Intuitively, all words in $\lang g$ share the same behaviour in $A$.
    We show that $\lang {g, h} \cap \lang A \neq \emptyset$ implies $\lang {g, h} \subseteq \lang A$.
    %
    To this end, consider an arbitrary word $w' \in \lang {g, h} = \lang g \cdot \lang h^\omega$
    which is naturally factorised as $w' = w'_0 w'_1 w'_2 \cdots$
    where $w'_0 \in \lang g$ and $w'_1, w'_2, \dots \in \lang h$.
    %
    By assumption, there exists a word $w \in \lang {g, h} = \lang g \cdot \lang h^\omega$
    which is accepted by the automaton $A$.
    %
    Thus let $\pi$ be an accepting run of $A$ on $w$.
    %
    We can write $w = w_0 w_1 w_2 \cdots$ where $w_0 \in \lang g$ and $w_1, w_2, \dots \in \lang h$.
    %
    We can split $\pi$ into segments
    %
    \begin{align*}
        \pi = p_0 \goesto {w_0} p_1 \goesto {w_1} p_2 \goesto {w_2} \cdots
    \end{align*}
    %
    \st~infinitely many segments visit an accepting transition.
    %
    Since $w_0, w_0'$ are both in $\lang g$, they share the same behaviour in $A$,
    and in particular $p_0 \goesto {w_0'} p_1$.
    %
    The same argument applies to $w_i, w_i' \in \lang h$,
    which means that $p_i \goesto {w_i'} p_{i+1}$ as well.
    %
    We thus obtain an infinite run
    %
    \begin{align*}
        \pi = p_0 \goesto {w_0'} p_1 \goesto {w_1'} p_2 \goesto {w_2'} \cdots
    \end{align*}
    %
    Moreover, the construction preserves accepting segments,
    and thus $\pi$ is an accepting run of $A$ on $w'$, as required.
\end{solution}

\begin{exercise}
    Show that the following problem is decidable, in polynomial time.
    %
    In input we are given a pair of graphs $g, h$
    and we need to decide whether $\lang {g, h} \cap \lang A \neq \emptyset$.
\end{exercise}
\begin{solution}
    Decompose $h$ into its strongly connected components (SCCs).
    Call a component \emph{accepting} if it contains an accepting transition.
    This can be done in polynomial time.
    %
    We have $\lang {g, h} \cap \lang A \neq \emptyset$
    iff there exists an edge $e = \tuple {p, x, q}$ of $g$ with $p \in I$ and $x \neq \bot$
    \st~in $h$ there is a path from $q$ to an accepting SCC of $h$.
    %
    This can also be checked in polynomial time.
\end{solution}

\begin{exercise}
    For every pair of graphs $g, h$,
    the language $\lang {g, h}$ is $\omega$-regular.
\end{exercise}
\begin{solution}
    First of all, the language of a single edge $\lang e \subseteq \Sigma^*$ is regular,
    as it easily follows from its definition.
    %
    The language of a graph is defined as a finite intersection of languages of edges, and is therefore regular.
    Finally, $\lang {g, h}$ is defined as a concatenation of a regular language $\lang g$
    and an $\omega$-power of a regular language $\lang h$, and is therefore $\omega$-regular.
\end{solution}

The exercises above show that we can write $\Sigma^\omega \setminus \lang A$
\emph{effectively} as a finite union of $\omega$-regular languages $\lang {g, h}$:
%
\begin{align}
    \Sigma^\omega \setminus \lang A
        = \bigcup_{g, h : \lang {g, h} \cap \lang A = \emptyset} \lang {g, h}.
\end{align}

\bibliographystyle{plainurl}
\bibliography{bibliography}

\end{document}